\section{Chương~\ref{sec:serocomputer}. Nơ-ron \nt{SERO}}

Các tính chất của chính nơ-ron \nt{SERO} (nhịp, tự ức chế, dấu theo thụ thể, các phân hệ) đã được đo trực tiếp; các phép tính của chương --- bộ điều chỉnh, bộ lọc, logic --- suy ra từ các phương trình của nó; $\tau_c$, hệ số khuếch tán của serotonin và số vị trí phóng là ước tính; vai trò bộ điều chỉnh mức của vòng lặp trong não là giả thuyết (trong nhân raphe, tự ức chế mang tính điểm và chỉ gây những quãng ngừng ngắn), còn vai trò tầm nhìn của \nt{SERO} là giả thuyết được các thí nghiệm hành vi ủng hộ.

{\footnotesize
\setlength{\tabcolsep}{3pt}\renewcommand{\arraystretch}{1.15}
\begin{longtable}{>{\raggedright}p{0.5cm}>{\raggedright}p{4.0cm}>{\raggedright}p{5.6cm}>{\raggedright}p{2.3cm}>{\raggedright\arraybackslash}p{1.7cm}}
\toprule
TT & Khẳng định & Thí nghiệm và điều đã đo & Nguồn & Trạng thái \\
\midrule\endfirsthead
\toprule
TT & Khẳng định & Thí nghiệm và điều đã đo & Nguồn & Trạng thái \\
\midrule\endhead
1 & Dấu tác động của serotonin do thụ thể của bên nhận chọn (mệnh đề~\ref{prop:receptor}) & Thụ thể serotonin được chia thành các họ theo dược lý và cơ chế: 5-HT$_{1A}$ qua protein G mở kênh kali và ức chế tế bào, 5-HT$_{2A}$ và thụ thể hướng ion 5-HT$_3$ gây kích thích. Cùng một chất dẫn truyền cho đáp ứng ngược nhau ở các tế bào khác nhau. & \cite{BarnesSharp1999} & đã đo \\
2 & Khi thức, nơ-ron \nt{SERO} phát đều vài xung mỗi giây & Ghi nơ-ron nhân raphe lưng ở mèo di chuyển tự do: phát xung nhịp nhàng chậm $2.8\pm0.2$ xung/s khi thức yên tĩnh; trong giấc ngủ sóng chậm tần số giảm $34$--$68\,\%$, trong giấc ngủ REM giảm $98\,\%$. Ở chuột cống gây mê: $1.7\pm0.5$ xung/s, rất đều. & \cite{TrulsonJacobs1979, GagnonParent2014, JacobsAzmitia1992} & đã đo \\
3 & Nơ-ron \nt{SERO} là bộ tạo nhịp: không cần được kích cho từng xung, nhưng cần dòng nền đều (trong lát cắt là noradrenaline qua thụ thể $\alpha_1$, trong cơ thể là từ \nt{NORA}) & Trong lát cắt não, tế bào phát xung chậm và đều, sau mỗi xung có một pha tăng phân cực sau xung dài $200$--$800\ms$. Trong lát cắt, số tế bào tự phát xung ít hơn trong cơ thể: nhịp đều cần đầu vào noradrenaline liên tục qua thụ thể $\alpha_1$, chặn chúng thì nhịp tắt. & \cite{VandermaelenAghajanian1983} & đã đo \\
4 & Hầu hết thụ thể là hướng chuyển hóa; đáp ứng chậm: tăng và giảm trong khoảng một giây & Mọi họ thụ thể, trừ 5-HT$_3$, đều gắn với protein G. Trong lát cắt, serotonin được phóng ra do kích thích tạo qua 5-HT$_{1A}$ một dòng ức chế tăng và giảm trong vòng một giây. & \cite{BarnesSharp1999, CourtneyFord2016} & đã đo \\
5 & Ở các vùng đích, serotonin thoát vào thể tích chứ không chỉ vào khe; ngay trong nhân raphe, ức chế lân cận đi theo kiểu điểm--điểm & Đo vôn-ampe vòng quét nhanh trong lát cắt: lượng phóng trên một xung như nhau ở vùng có khớp thần kinh và ở nhân raphe, nơi gần như không có khớp thần kinh; không phụ thuộc vào việc chặn tái hấp thu và thụ thể; số phân tử trên mỗi đầu tận cùng ít hơn nhiều so với số chất vận chuyển và thụ thể, còn nồng độ ngoại bào trùng với ái lực của thụ thể chính. Trong nhân raphe, ức chế qua 5-HT$_{1A}$ đi theo kiểu điểm--điểm: chất vận chuyển không để serotonin tích tụ ngoài khe. & \cite{Bunin1998, CourtneyFord2016} & đã đo \\
6 & Nồng độ theo~\eqref{eq:seroneuron} giảm do bị bơm đi; $\tau_c$ cỡ một giây là ước tính & Đo vôn-ampe trong não sống: serotonin ngoại bào bị giới hạn trước hết bởi tái hấp thu và phân hủy, chứ không phải bởi tổng hợp. Các công trình được trích không đo trực tiếp $\tau_c$; bậc độ lớn là ước tính của chương. & \cite{Hashemi2012serotonin} & đã đo; $\tau_c$ là ước tính \\
7 & NOT và NOR trên một bộ tạo nhịp; tính đầy đủ của logic \nt{SERO} (mệnh đề~\ref{prop:serocomplete}, hình~\ref{fig:serogates}) & Suy ra từ~\eqref{eq:seroneuron}: bộ tạo nhịp bị ức chế là NOR, và NOR đầy đủ về chức năng. Chưa có thí nghiệm nào lắp logic từ nơ-ron \nt{SERO}; điều kiện các thể tích tách biệt không được thỏa mãn trong não. & suy ra trong chương & lý thuyết \\
8 & Serotonin ức chế chính nơ-ron \nt{SERO} qua các thụ thể trên nó & Ở chuột cống gây mê, chất chủ vận 5-HT$_{1A}$ tiêm tĩnh mạch và điện di ion ức chế sự phát xung của nơ-ron raphe lưng với cùng quan hệ liều--đáp ứng như chính serotonin; chất chủ vận 5-HT$_{1B}$ gần như không tác dụng. Trong lát cắt, serotonin nội sinh từ các tế bào lân cận tạo qua 5-HT$_{1A}$ một dòng và một quãng ngừng ngắn trong chuỗi xung, không làm đổi tần số trung bình. & \cite{Sprouse1987, CourtneyFord2016} & đã đo \\
9 & Trong mô hình, vòng lặp qua thể tích chung là bộ điều chỉnh mức: sai lệch và hằng số thời gian giảm $1+G$ lần~\eqref{eq:feedback}; việc vòng lặp hoạt động như vậy trong não là giả thuyết & Công thức kinh điển của bộ khuếch đại có hồi tiếp âm; được suy lại trong chương. Hệ số khuếch đại vòng $G$ của hệ \nt{SERO} chưa được đo. Đã đo: trong lát cắt, tự ức chế đi theo kiểu điểm--điểm, gây quãng ngừng ngắn và không đổi tần số trung bình. & \cite{Black1934feedback, CourtneyFord2016} & lý thuyết; trong não là giả thuyết \\
10 & Nồng độ là trung bình trượt của tần số, bộ lọc thông thấp~\eqref{eq:lowpass}, hình~\ref{fig:lowpass} & Nghiệm của phương trình tuyến tính~\eqref{eq:seroneuron}; hình vẽ được tính theo công thức này với $\tau_c=1\,\text{s}$. Không có mô hình riêng trong \texttt{models/}. & suy ra trong chương & lý thuyết \\
11 & Độ ngoằn ngoèo của khe làm khuếch tán của phân tử nhỏ chậm đi khoảng $2.6$ lần & Phương pháp nguồn điểm: điện di ion tetramethylammonium và ghi nồng độ của nó bằng điện cực chọn lọc ion trong lát cắt và trong não sống. Độ ngoằn ngoèo $\lambda\approx1.6$, tức $D$ hiệu dụng nhỏ hơn giá trị tự do $\lambda^2\approx2.6$ lần; tỷ lệ thể tích gian bào khoảng $0.2$. Hệ số khuếch tán của chính serotonin trong mô không được đo trong các công trình này; trong chương nó là ước tính. & \cite{Sykova2008} & đã đo \\
12 & Chiều dài của một ``dây dẫn'' độc lập $\ell\sim\sqrt{D\tau}\approx20\um$~\eqref{eq:diffusionlength}; số dây bị giới hạn bởi số vùng như vậy (mệnh đề~\ref{prop:volumewires}) & Định luật khuếch tán và thế các ước tính $D$, $\tau$ (dòng~6 và~11); mệnh đề~\ref{prop:volumewires} là hệ quả. Chưa có thí nghiệm trực tiếp đếm số kênh độc lập của truyền dẫn thể tích. & suy ra trong chương & lý thuyết \\
13 & Đầu ra của một nơ-ron \nt{SERO} trải trên những vùng não rất lớn; số vị trí phóng ước tính $\sim10^5$ & Tái dựng toàn bộ từng nơ-ron raphe lưng ở chuột cống: mọi sợi trục đi lên đều phân nhánh mạnh, tổng chiều dài sợi trục của một tế bào tới $18.7$~cm, các nhánh của một tế bào vươn tới thể vân, vỏ não trước trán và hạch hạnh nhân. Số vị trí phóng trên mỗi tế bào chưa được đếm trực tiếp. & \cite{GagnonParent2014} & đã đo; $10^5$ là ước tính \\
14 & Các nơ-ron \nt{SERO} chia thành vài phân hệ với đầu ra và đáp ứng khác nhau & Đánh dấu di truyền bằng virus ở chuột nhắt: các nơ-ron chiếu tới hạch hạnh nhân và tới vỏ não trán hầu như không chồng lấn về phân nhánh, nhận đầu vào khác nhau và đáp ứng ngược nhau với kích thích khó chịu; bật và tắt mỗi nhóm làm thay đổi hành vi theo cách khác nhau. & \cite{Ren2018} & đã đo \\
15 & Điều kiện kiên nhẫn $D<\tau_\gamma\ln(R/r_0)$~\eqref{eq:patience} với chiết khấu hàm mũ; với chiết khấu đo được, dạng hyperbol, là $D<\tau_\gamma(R/r_0-1)$ & Suy ra từ chiết khấu $e^{-D/\tau_\gamma}$ hoặc $1/(1+D/\tau_\gamma)$; suy ra trong chương. Khỉ, lựa chọn với độ trễ cỡ giây: chiết khấu gần hyperbol hơn hàm mũ; đáp ứng của nơ-ron \nt{DOPA} với tín hiệu báo trước phần thưởng hoãn lại cũng giảm theo độ trễ như hyperbol. & suy ra trong chương; \cite{KobayashiSchultz2008} & lý thuyết \\
16 & \nt{SERO} đặt tầm nhìn $\tau_\gamma$ (mục~\ref{sec:horizon}) & Kích hoạt quang di truyền các nơ-ron \nt{SERO} raphe lưng ở chuột nhắt kéo dài thời gian chờ phần thưởng hoãn lại và tăng sự kiên trì khi tìm phần thưởng đã trở nên hiếm. Ở người, giảm serotonin (chế độ ăn không có tryptophan) làm chiết khấu phần thưởng hoãn lại dốc hơn. Nồng độ chuyển thành $\tau_\gamma$ thế nào thì chưa rõ. & \cite{Doya2002, Miyazaki2014, Lottem2018, Schweighofer2008} & giả thuyết \\
\bottomrule
\end{longtable}}
